\section{问题重述}

\subsection{问题背景}
% TODO: 用 2-3 句概括问题背景（对象、场景、核心矛盾）。禁止套话（"数学建模是解决实际问题的
% 重要工具"等）与虚构业务事实；背景必须来自题目原文。

\subsection{问题任务}
% TODO: 按 QUESTION_CONTRACT.questions 逐条列任务（共 \questioncount 问）。
% 每问给出：输入（数据/变量）、输出（要求的结果形态）、核心约束。
\questiontask{问题一}{% TODO: 输入/输出/约束
}

\subsection{各问目标与输出}
% TODO: 只用竞赛语言概括每问要求给出的数值、方案、时点、区间或判定方法；评价指标写成
% 数学/统计定义。不得写 gate、registry、authority、hash、closure、验收证据等后台机制。
% 仅当题目确实需要单独汇总输出时保留本小节；纯分析题可删。
